\chapter{Discussion}
\label{ch:discussion}

\section{Interpreting the inverted premium: commitment without reception}
\label{sec:disc-premium}

The headline result is a reversal: without communication the strong models out-earn the
weak, and with it they fall behind. The natural reading of that sentence is the wrong one.
``Communication'' suggests the weak agents heard the strong agents' cooperative promises and
took advantage of them, but that is not what the apparatus does. As
Section~\ref{sec:communication} sets out, a message is produced and logged and never
delivered; no agent in any game reads another's words. Whatever inverts the premium does so
without a single pledge changing hands.

What the communication condition actually adds to a round is one thing: the request that an
agent state a public intention before it acts. The behavior that follows says the request is
enough. Asked to commit, the capable models commit, and then they keep the commitment. The
transcript in Section~\ref{sec:res-premium} is typical: a strong agent announces it will give
its whole endowment and does so in every round, while the weak agents around it contribute a
fraction and pocket the rest. The strong agent is not exploited because anyone read its
promise; it is exploited because it cooperates, visibly, and the visible contribution is what
the others free-ride on. The history they all see carries contributions, and a steady high
contribution is a standing invitation to give less.

The asymmetry, then, is in who honors a stated intention. Both tiers are asked to pledge, and
both pledge cooperatively: the classifier of Section~\ref{sec:res-deception} labels nearly
nine in ten strong messages as cooperative promises. The difference is the follow-through. The
capable models act on the commitment they articulated; the weak models say the same words and
then do otherwise, which is exactly what their higher surprisal records. A public commitment
moves the one who makes it, and moves the capable more than the weak, because only the capable
treat their own stated policy as binding. That is why the effect needs no audience: the
commitment does its work on the committer.

This is a narrower and stranger claim than the word ``communication'' invites, and its size is
worth being exact about. The study does not show that talking helps weak agents exploit strong
ones, nor that cheap talk coordinates a group. It shows that asking a capable agent to say what
it will do makes it more likely to do the cooperative thing it said, at a cost, in a setting
where cooperation can be taken advantage of. This sits beside a recent result that runs the
other way. \emph{More Capable, Less Cooperative}~\cite{morecapable2026} finds that where helping
is free, more capable models can cooperate \emph{less}, withholding help they would lose nothing
by giving; here the capable models cooperate \emph{more}, and are punished for it. The two are not
in tension but two faces of one decoupling of capability from cooperative advantage, in one case
failing to produce cooperation that costs nothing, in the other producing cooperation that is
exploited. What this study adds is the lever, a public pledge, and the setting, a group dilemma,
in which the second face appears.

\section{Why the dissociation matters: group vs. dyadic structure}
\label{sec:disc-dissociation}

The inversion is a fact about the Public Goods Game. In the Trust Game the same populations
produce a premium that sits near zero and never separates cleanly from it
(Section~\ref{sec:res-dissociation}). Read carelessly, that looks like a failure to replicate;
read for what it is, it locates the effect. The two games differ in one structural way that
matters here, and the difference is why the lever pulls in one and not the other.

The Public Goods Game has a crowd. Four players draw on a shared pool, and one agent's
contribution is visible to, and divisible among, the other three. A capable agent that commits
to contributing creates exactly what free-riding needs: a standing benefit others can take a
share of without paying in. The commitment and the crowd are both required. Take away the
commitment and the strong agent contributes no more than anyone else; take away the crowd and
no one is positioned to free-ride on its contribution. The Trust Game removes the crowd. It is a
dyad, and its second mover faces a direct choice of how much to return, not a chance to draw
quietly from a pool someone else filled. A capable trustee can keep the tripled stake, but doing
so is a visible refusal to one partner, not an anonymous skim from a group.

The dissociation is the cross-game design earning its place. Had the study run only the Public
Goods Game, the inverted premium could have read as a general property of capable agents in
mixed company. Running the Trust Game beside it shows that it is not: the effect needs group
structure, a shared resource and an audience of beneficiaries, and fades where that structure is
absent. That is more useful than a uniform effect would have been, because it says where to
expect the behavior and where not to. The pre-registered test for H5 asked whether the effects
reverse between the games; they do not, so the prediction is not falsified by its own criterion
(Section~\ref{sec:res-dissociation}). The stronger reading is the dissociation itself: the
capability premium and its inversion are phenomena of the group game.

\section{A methodological contribution: a valid vs. an artefactual deception measure}
\label{sec:disc-deception}

The deception result is also a methodological one, because two measures built on the same idea
disagree, and the disagreement is the point. Both ask the stated-versus-revealed question a
forced-choice study put to single decisions~\cite{alignmentrevisited2025}: does an agent's
action match the policy its words implied? Lifted to a repeated game, the question splits into
two measures that need not agree, and here they point opposite ways.

The contemporaneous surprisal pairs each message with the action of the same round and asks how
much weight the stated policy gave to the move actually made. By it, the weak tier is the
deceptive one, and the blind classifier, which never sees an action, reaches the same verdict
from the text alone (Section~\ref{sec:res-deception}). The global policy-KL instead compares an
agent's whole-game distribution of actions against its stated policies, and it calls the strong
tier the deceptive one, which is backwards. The reason is concentration, not deceit. A steady
cooperator has a sharp action distribution, and a sharp distribution diverges far from any hedged
statement of intent, so the measure scores consistency as dishonesty: the most reliable agents
come out looking the least honest.

This is why the thesis keeps the global divergence only as a negative control and reports the
contemporaneous surprisal as the signal. The contribution is less a new metric than a
demonstration that the obvious one fails in a specific, instructive way. A stated-versus-revealed
gap measured over pooled behavior confounds deception with regularity; only the action-by-action
version, scored where a word meets the deed it was meant to predict, separates the two. A measure
that called the steadiest cooperators the worst liars would have quietly produced the wrong
story, and reporting it beside the one that works, rather than dropping it, is what lets the
deception finding rest on the agreement of methods rather than the luck of a choice.

\section{Implications for deploying mixed agent systems}
\label{sec:disc-implications}

If capable and weak language models are going to share tasks, and increasingly they will, the
result carries a few practical warnings. The first is that capability does not buy advantage in a
shared-resource setting. A stronger model dropped into a group that draws on a common pool can
end up subsidizing the others, and the asymmetry is largest exactly where the design might seem
safest, when the capable agent is asked to be transparent about what it means to do. Asking an
agent to declare its plan is a common move in alignment and oversight; in a group dilemma it is
the move that costs the capable agent most, because it turns a private disposition to cooperate
into a public, exploitable commitment.

The second warning is about measuring honesty in these systems. A deployer who wants to flag
deceptive agents has an obvious tool, the divergence between what an agent says and what it does,
and the obvious form of that tool is wrong in a way that matters: pooled over a run, it fingers
the most consistent cooperators and misses the agents whose words and deeds part company turn by
turn. The measure has to be contemporaneous to mean anything. The third is a scope condition more
than a warning: the effect is a property of group structure, not of capability as such, so a
one-on-one delegation between a strong and a weak model is a different situation from a many-agent
commons and should not be expected to behave the same way.

None of this is a deployment study, and the distance between four language models in a tokenized
game and a working multi-agent system is wide. What the results offer is not a prediction about
production but a caution about intuitions: that the strong agent will dominate, that letting
agents talk can only help, and that any reasonable measure of stated-versus-revealed behavior
captures deception. All three are natural, and all three are wrong in the settings studied here.

\section{Limitations}
\label{sec:disc-limitations}

The honest reading of these results depends on holding several limitations in view at once, most
recorded already where they arose and gathered here so none is lost.

The first concerns the communication condition, and it is the most important. What the study
calls communication is not an exchange. Each agent is asked to make a public pledge and is told
the others will see it, but the runner never delivers the message
(Section~\ref{sec:communication}); agents act on the contribution history, not on each other's
words. The condition is public commitment without reception, and every claim about it here is
meant in those terms. The word overstates the apparatus, and the finding is the narrower one of
Section~\ref{sec:disc-premium}, not a result about agents reading and answering one another.

The deception measures carry their own caveats, declared with the metric. The model that scores
an agent's stated policy is Claude Haiku~4.5, which also plays as one of the strong agents, so
the scorer is in part one of the scored (Section~\ref{sec:integrity}); the coupling is mitigated
by reading the same conclusion off three methods that share no machinery, the contemporaneous
surprisal, a blind text classifier, and the global divergence that runs the other way as a
control, rather than off any single read. The surprisal was validated against human labels by a
single rater, the metric's author, working blind and single-pass against a fixed rubric, and it
cleared significance but not the pre-registered correlation bar, which is why it is reported as
exploratory rather than confirmatory. These are real limits on how much the deception result can
bear; it is offered as a methodological demonstration and a convergence of signals, not a
calibrated measurement.

The capability axis is coarse. ``Strong'' means three specific frontier models and ``weak'' means
one small local one, and the tiers differ in more than a single dimension of capability, so the
results speak to that contrast rather than to capability as a clean variable. The scale is modest
too: 320 games, with the per-game deception divergence sparse enough that the signal lives in the
spread across games rather than in any one of them. The primary regression keeps only a per-game
random intercept where the pre-registration named one for agents as well, because the modelled
outcome is a single end-of-game payoff per agent (Section~\ref{sec:res-premium}); the models'
outputs are not bit-reproducible regardless of seed; and the network reading of
Section~\ref{sec:res-networks} draws on the same behavioral features as the tier classifier, so
the two are not independent evidence. None of these sinks the central results, which are large
and consistent, but each marks a place where the reader should not ask the data for more than it
holds.

\section{Future work}
\label{sec:disc-future}

The clearest next step follows from the mechanism. If a public commitment changes behavior
without being received, the experiment to run is to vary reception directly: deliver the messages
in one condition and withhold them in another, holding the pledge-request fixed, and see whether
being read changes what a commitment does. The present design cannot separate the request to
pledge from the claim that others will see it, because it never delivers; a channel that actually
carried messages would let the framing effect and any genuine exchange effect be told apart, and
would test whether the inversion survives once the audience is real.

Two other extensions would firm up what the results can claim. The capability axis wants
widening, from three strong models against one weak to a graded set that varies capability while
holding other differences as nearly fixed as possible, so that ``capability'' becomes something
closer to a measured variable than a contrast between a frontier API and a small local model. And
the deception measure wants the validation the pilot only sketched: more messages, more than one
rater, and a labeling effort large enough to say whether the contemporaneous surprisal tracks
human judgments well enough to stand on its own rather than only in the company of other signals.
Beyond these, the dissociation invites a map. Varying group size and the structure of the shared
resource would trace the boundary that the Public Goods and Trust games mark at only two points,
and show how much of a crowd a commitment needs before it can be exploited.
